\chapter{Introduction}
\label{ch:ch1_intro}

% [DRAFT v1] Source material: docs/plan.md, docs/arquitectura.md, docs/via_asic.md.

\section{Motivation: keys are only as good as their entropy}
\label{sec:motivation}

Every secure connection an embedded device opens, every firmware image it
authenticates and every credential it stores rests on a key, and every
key rests on a source of randomness. If that source is weaker than
assumed, no amount of cryptographic sophistication downstream recovers
what was lost at the start: a 128-bit key drawn from a generator with
effectively 40 bits of entropy is a 40-bit key. This is not a theoretical
concern. Frequency injection on the supply rail of a commercial secure
microcontroller collapsed its key space from $2^{64}$ to about
3300~\cite{markettos_2009_frequency}; a certified smart-card generator was
found biased despite passing the standard statistical
suites~\cite{hurleysmith_2018_certifiably}; and the hardware generator of
one of the most widely deployed IoT microcontrollers, the ESP32, silently
returns pseudo-random values when its radio is off while still passing
every black-box test~\cite{blancoromero_2026_entropy}.

The common thread is that randomness has been judged by looking at the
output. Statistical batteries such as NIST SP~800-22 check whether a bit
stream \emph{looks} uniform and independent; they cannot tell a physical
noise source from a deterministic generator seeded with a constant, and
NIST said so formally in 2022 when it announced the revision of that
publication, ``rejecting its use for assessing cryptographic random
number generators''~\cite{nist_2022_decision}. The current standards
--- AIS~20/31 in Europe~\cite{peter_2024_ais31} and SP~800-90B in the
United States~\cite{turan_2018_sp80090b} --- ask for something different:
a stochastic model of the noise source, an estimate of its min-entropy
derived from that model, and health tests that watch the source itself
rather than its output. That is a microelectronics problem before it is a
cryptography problem, and it is the problem this thesis addresses.

\section{From the black-box ESP32 study to a white-box engine}
\label{sec:background}

This work continues a previous project by the author that characterized
the hardware random number generator of the ESP32 as a black box: the
same peripheral was exercised under four operating conditions, its output
was run through a self-validated implementation of the fifteen tests of
SP~800-22, and a key-harvesting scheme combining several on-chip sources
was proposed. Its central finding was methodological: the tests passed
under conditions in which the entropy could not have been what the tests
implied, which is exactly the blind spot described above. Two of its
numbers also had to be revisited with hindsight --- a temperature
correlation was reported from a sensor whose resolution was a single
step over the observed range, and a stability figure was inferred where
it could have been measured --- and both lessons are turned into rules in
Section~\ref{sec:conventions}.

The natural next step was to build the generator instead of measuring
someone else's: an engine whose entropy source is designed, whose jitter
is measured inside the device, whose min-entropy follows from a model,
and whose keys are produced and used by primitives verified against the
standards. The ESP32 changes role from subject to judge. It has its own
AES and SHA accelerators and a mature cryptographic library, so it can
recompute everything the engine produces with an implementation that
shares nothing with it, over a real bus.

The engine is designed for silicon. Secure elements --- small integrated
circuits that generate, store and use keys on behalf of a host
microcontroller over I2C --- are a mature product category; the
\emph{transparency} of their entropy is not. Of three widely used parts,
one ships with a NIST validation of its entropy source and two do not
mention SP~800-90 anywhere in their documentation~\cite{microchip_atecc608b_ds}.
The FPGA in this work is therefore the verification vehicle, not the
target: it lets millions of test vectors run against the design in
minutes, which no transistor-level simulator can, before the same
description is synthesized on standard cells.

\section{Objectives}
\label{sec:objectives}

\begin{enumerate}
\item[O1] Design a ring-oscillator entropy source at transistor level,
with enable and sampler, and simulate it with noise, stating the
numerical resolution of the bench before any figure is read from it.
\item[O2] Extract the jitter from the simulated waveforms with an
estimator that separates quantization, thermal and flicker components,
validated beforehand on synthetic data of known jitter.
\item[O3] Turn the jitter into a min-entropy bound through a stochastic
model checked against an exact computation, and derive the sampler
divider that meets the AIS~20/31 target with margin.
\item[O4] Write the digital engine --- health tests, AES-based
conditioning and CTR\_DRBG, AES-CMAC and an I2C register interface --- as
technology-independent RTL verified against official vectors.
\item[O5] Verify the engine functionally on an FPGA with the ESP32 as an
independent oracle: known-answer tests, cross interoperability and
challenge-response.
\item[O6] Run an entropy campaign on real silicon (the FPGA's rings) with
the SP~800-90B estimators and compare with the ESP32 generator under the
same methodology.
\item[O7] Synthesize the engine on an open standard-cell library and
report its area in gate equivalents with the growth factors to core area.
\item[O8] Compare the design with the ESP32 generator, with commercial
secure elements and with published circuits, with stated uncertainties.
\end{enumerate}

The minimum defensible result is O1--O3 together with O5: an entropy
source designed from the transistor, with its jitter extracted and turned
into a bound, and a digital engine verified against an independent
implementation. A negative result that is well measured closes the work
as well as a positive one: if the accumulation exponent on hardware is
two rather than one, flicker dominates and the divider cannot grow
indefinitely, and that is a finding.

\section{Scope and exclusions}
\label{sec:scope}

Four things are deliberately out of scope, each for a stated reason.
\emph{Fabrication}: an academic shuttle is affordable and the engine fits
in the smallest block offered, but the schedule from tape-out to packaged
and characterized samples exceeds a master's thesis; the cost and the path
are documented, not executed. \emph{Asymmetric cryptography}: an
elliptic-curve multiplier is a project in itself and would consume the
whole area budget; the engine provides symmetric keys and authentication.
\emph{Side-channel resistance}: the threat model is a logical, remote
adversary; timing leaks are avoided by construction, but no masking or
trace-based evaluation is attempted, and the institute's own laboratory
for such work is noted as the natural continuation. \emph{Formal
certification}: the criteria of AIS~20/31 and SP~800-90B and the official
vectors are used throughout, and no validation is claimed.

\section{Methodology and reproducibility conventions}
\label{sec:conventions}

Four conventions apply to every figure in this document. First, each
number is labelled \meas{}, \simu{} or \est{}, and carries its
uncertainty or the source it comes from. Second, \textbf{the resolution of
the instrument is declared before any trend is interpreted}, whether the
instrument is a sensor, an estimator or a simulator --- the numerical
noise floor of the SPICE bench and the quantization of a temperature
sensor are treated alike. Third, quantities are measured rather than
inferred whenever a direct measurement exists. Fourth, novelty is stated
as ``not found in the literature consulted'', never as ``does not exist'';
one such claim in an early draft of this work had to be corrected when a
2024 paper measuring exactly the quantity in question was found, and the
correction is left visible in the working notes.

Everything is released: the VHDL and its testbenches, the Python models
and analysis, the firmware, the synthesis flow, the official test vectors
used and the raw data of every campaign, with the commands that reproduce
each table.

\section{Structure of the document}
\label{sec:structure}

Chapter~\ref{ch:ch2_soa} reviews ring-oscillator generators, their
stochastic models, the standards and the commercial landscape.
Chapter~\ref{ch:ch3_entropy} builds the chain from transistor to
min-entropy and validates every link. Chapter~\ref{ch:ch4_engine}
describes the digital engine and Chapter~\ref{ch:ch5_verification} its
verification. Chapter~\ref{ch:ch6_asic} synthesizes the engine on standard
cells and states the path and cost of an integrated circuit.
Chapter~\ref{ch:ch7_results} reports the hardware campaign and the
comparisons, and Chapter~\ref{ch:ch8_concl} concludes.
